\documentclass[11pt]{article}
\usepackage[margin=1in]{geometry}
\usepackage[T1]{fontenc}
\usepackage{lmodern}
\usepackage{microtype}
\usepackage{xcolor}
\usepackage{hyperref}
\usepackage{booktabs}
\usepackage{enumitem}
\usepackage{listings}
\usepackage{tcolorbox}
\usepackage{parskip}
\usepackage{array}
\newcolumntype{L}[1]{>{\raggedright\arraybackslash}p{#1}}

% \path{} (from url, loaded by hyperref) prints paths and identifiers
% verbatim -- so `--` is not turned into a dash -- and may break at / _ - .
\makeatletter
\g@addto@macro\UrlBreaks{\do\_\do\-\do\.}
\makeatother
\urlstyle{tt}
\setlength{\emergencystretch}{3em}

\hypersetup{
  colorlinks=true,
  linkcolor=blue!60!black,
  urlcolor=blue!60!black
}

\definecolor{codebg}{RGB}{245,245,245}
\definecolor{codecomment}{RGB}{90,100,90}
\definecolor{codekeyword}{RGB}{0,70,140}
\definecolor{codestring}{RGB}{150,60,20}

\lstdefinestyle{pythonstyle}{
  language=Python,
  backgroundcolor=\color{codebg},
  basicstyle=\ttfamily\small,
  keywordstyle=\color{codekeyword}\bfseries,
  commentstyle=\color{codecomment}\itshape,
  stringstyle=\color{codestring},
  showstringspaces=false,
  breaklines=true,
  frame=single,
  rulecolor=\color{black!20},
  columns=fullflexible,
  keepspaces=true
}
\lstdefinestyle{shellstyle}{
  language=bash,
  backgroundcolor=\color{codebg},
  basicstyle=\ttfamily\small,
  showstringspaces=false,
  breaklines=true,
  frame=single,
  rulecolor=\color{black!20},
  columns=fullflexible,
  keepspaces=true
}
\lstset{style=pythonstyle}

\newtcolorbox{keyidea}{
  colback=blue!5, colframe=blue!50!black,
  title=Key Idea, fonttitle=\bfseries
}
\newtcolorbox{tipbox}{
  colback=green!5, colframe=green!45!black,
  title=Tip, fonttitle=\bfseries
}
\newtcolorbox{warningbox}{
  colback=orange!7, colframe=orange!60!black,
  title=Common Pitfall, fonttitle=\bfseries
}

\title{
  \textbf{Codebase Change Summary --- Week 6}\\[4pt]
  \large Data-Flow Testing and Coverage Concepts\\[2pt]
  \normalsize SE 413/513 --- Software Testing
}
\author{Demo Testing Suite}
\date{}

\begin{document}
\maketitle

\begin{center}
\textit{What changed in \path{demo-testing-suite/} since Week 5, and why}
\end{center}

\vspace{0.3cm}

% ============================================================
\section{Purpose of This Document}
% ============================================================
The \path{demo-testing-suite/} codebase is a single, shared project that
grows one lecture at a time rather than being rebuilt from scratch each
week. This document summarizes exactly what was added or changed for
Week 6 and, more importantly, explains \emph{why} each change exists ---
which concept from the Week 6 reading, \textit{Data-Flow Testing and
Coverage Concepts}, it demonstrates. Read this alongside the reading and
\path{demo-testing-suite/README.md}.

\begin{keyidea}
Nothing from Weeks 1--5 was removed or rewritten. Everything below is
\emph{additive}: the running examples you already saw are all still
there, unchanged, including their two intentional \texttt{xfail} results
and one intentional \texttt{FAIL}.
\end{keyidea}

% ============================================================
\section{What Changed At a Glance}
% ============================================================
\begin{center}
\small
\begin{tabular}{@{}L{5.9cm}L{1.5cm}L{7.0cm}@{}}
\toprule
\textbf{File} & \textbf{Status} & \textbf{Purpose} \\
\midrule
\multicolumn{3}{@{}l}{\textit{The main suite}}\\
\path{app/billing/charges.py} & New & The reading's worked example,
\path{sum_positive()}, verbatim with its S1--S5 labels: sums the
charges on a statement, skipping credits. \\
\path{tests/billing/test_charges.py} & New & Four tests whose ids name
the coverage criterion and the DU pairs each one covers, including the
loop-carried pair. \\
\midrule
\multicolumn{3}{@{}l}{\textit{Deliberately buggy illustrations (not in the default run)}}\\
\path{demos/data_flow/accumulator_bug.py} & New & \path{sum_positive()}
with the loop-carried pair broken. 100\% branch coverage; one All-Uses
test fails. \\
\path{demos/data_flow/receipt.py} & New & A correct function where a
one-test All-Defs set leaves a branch untested (75\%). \\
\path{demos/data_flow/anomalies.py} & New & A use-before-def and a dead
definition that tests pass over and \texttt{pylint} finds statically. \\
\midrule
\multicolumn{3}{@{}l}{\textit{Tooling and documentation}}\\
\texttt{requirements.txt} & Modified & Adds \texttt{pylint}. \\
\path{docs/DEMO_GUIDE_week06_data_flow.md} & New &
Section-by-section lecture companion for the reading. \\
\texttt{README.md} (suite and repository) & Modified & Updated layout,
\texttt{pylint} command, expected output, and the Weekly Updates index. \\
\bottomrule
\end{tabular}
\end{center}

\begin{tipbox}
Install the new dependency before running anything this week:
\begin{lstlisting}[style=shellstyle]
cd demo-testing-suite
pip install -r requirements.txt
\end{lstlisting}
\end{tipbox}

% ============================================================
\section{The Worked Example: \texttt{sum\_positive()}}
% ============================================================
\textbf{Locations:} \path{app/billing/charges.py},
\path{tests/billing/test_charges.py}

\begin{lstlisting}
def sum_positive(numbers):
    total = 0                  # S1
    for n in numbers:          # S2
        if n > 0:              # S3
            total = total + n  # S4
    return total               # S5
\end{lstlisting}

The module docstring repeats the reading's def/use table. The test file's
docstring repeats its seven DU pairs, with a ``covered by'' column. Each
test id says which criterion the test belongs to:

\begin{center}
\footnotesize
\begin{tabular}{@{}lllL{5.6cm}@{}}
\toprule
\textbf{Test id} & \textbf{Input} & \textbf{Result} & \textbf{New DU pairs (reading's numbering)} \\
\midrule
\texttt{all-defs:credit-only} & \texttt{[-1]} & \texttt{0} & \#2, \#6 \\
\texttt{all-defs:one-charge} & \texttt{[3]} & \texttt{3} & \#1, \#4, \#5, \#7 \\
\texttt{all-uses:two-charges} & \texttt{[3, 5]} & \texttt{8} & \#3 (loop-carried) \\
\path{du-path:loop-carried-via-credit} & \texttt{[3, -2, 5]} & \texttt{8} & none new: \#3 along a second def-clear path \\
\bottomrule
\end{tabular}
\end{center}

\begin{lstlisting}[style=shellstyle]
pytest tests/billing/test_charges.py -v \
    --cov=app.billing.charges --cov-branch --cov-report=term-missing
\end{lstlisting}

The first two tests already give 100\% branch coverage of this function.
The third adds no statement and no branch --- it adds a \emph{pair}. The
fourth reaches pair \#3 through the other back edge (S3's false outcome,
a skipped credit between two charges). It covers no new DU pair, so
All-Uses does not require it; it is the kind of extra path All-DU-Paths
asks for.

% ============================================================
\section{Why the Loop-Carried Pair Matters}
% ============================================================
\textbf{Locations:} \path{demos/data_flow/accumulator_bug.py},
\path{demos/data_flow/test_accumulator_bug.py}

The same function, with S4 broken:
\begin{lstlisting}
            total = 0 + n      # S4  KNOWN BUG -- should be total + n
\end{lstlisting}
Each charge now overwrites (\emph{kills}) the running total instead of
adding to it, and pair \#3 no longer exists. The demo's tests are the
All-Defs set plus one All-Uses test:
\begin{lstlisting}[style=shellstyle]
pytest demos/data_flow/test_accumulator_bug.py -v \
    --cov=accumulator_bug --cov-branch --cov-report=term-missing
\end{lstlisting}
\begin{lstlisting}[style=shellstyle]
Name                                 Stmts   Miss Branch BrPart  Cover   Missing
--------------------------------------------------------------------------------
demos/data_flow/accumulator_bug.py       6      0      4      0   100%
\end{lstlisting}
The two All-Defs tests pass, with \textbf{100\% statement and branch
coverage}. Only \path{test_all_uses_two_charges_exposes_the_bug} fails,
with \texttt{assert 5 == 8}: the function returns the last charge, not
the sum.

\begin{warningbox}
Branch coverage confirms every decision ran both ways. It says nothing
about whether an accumulated value was ever checked after more than one
update. Whenever a variable is both defined and used inside a loop, ask
whether the definition also reaches a use on a \emph{later} iteration.
\end{warningbox}

% ============================================================
\section{All-Defs Does Not Subsume Branch Coverage}
% ============================================================
\textbf{Locations:} \path{demos/data_flow/receipt.py},
\path{demos/data_flow/test_receipt.py}

The reading states this but its own example happens not to show it:
in \path{sum_positive()}, the All-Defs set also covers both branches.
\path{total_with_tax(amount, show_breakdown)} does show it. A single
call with \path{show_breakdown=False} gives every definition one use,
which is All-Defs:
\begin{lstlisting}[style=shellstyle]
pytest demos/data_flow/test_receipt.py \
    --cov=receipt --cov-branch --cov-report=term-missing
\end{lstlisting}
\begin{lstlisting}[style=shellstyle]
Name                         Stmts   Miss Branch BrPart  Cover   Missing
------------------------------------------------------------------------
demos/data_flow/receipt.py       6      1      2      1    75%   33
\end{lstlisting}
Line 33, the \texttt{print()} on the \texttt{True} branch, never ran.
Uncomment \path{test_total_with_breakdown} at the bottom of the test
file: it covers the remaining DU pairs (All-Uses) and the report
reaches 100\%.

% ============================================================
\section{Anomalies That Static Analysis Finds}
% ============================================================
\textbf{Locations:} \path{demos/data_flow/anomalies.py},
\path{demos/data_flow/test_anomalies.py}

Two of the reading's defect patterns, shipped on purpose:
\begin{itemize}
\item \textbf{Use before def:} \path{describe_balance()} never assigns
\texttt{status} when the balance is exactly 0.
\item \textbf{Dead definition:} \path{apply_payment()} computes
\texttt{remaining} and then returns the wrong variable.
\end{itemize}
The three tests pass, and every statement runs. \texttt{pylint} needs no
tests at all:
\begin{lstlisting}[style=shellstyle]
pylint demos/data_flow/anomalies.py --disable=all \
    --enable=possibly-used-before-assignment,unused-variable
\end{lstlisting}
\begin{lstlisting}[style=shellstyle]
anomalies.py:36:25: E0606: Possibly using variable 'status' before assignment
anomalies.py:40:4: W0612: Unused variable 'remaining'
\end{lstlisting}
This is the reading's ``Data-Flow Testing in Practice'' section:
data-flow \emph{coverage} tools are rare, but def-use \emph{analysis}
runs inside everyday linters, compilers, and IDEs.

\begin{tipbox}
In Python, a use before def crashes with \texttt{UnboundLocalError}
(uncomment \path{test_zero_balance} to see it). In some other languages
the same mistake silently returns garbage or a stale value, which is why
the reading treats it as a data-flow defect, not just a crash.
\end{tipbox}

% ============================================================
\section{Updated: Expected Test Results}
% ============================================================
From the project root:
\begin{lstlisting}[style=shellstyle]
cd demo-testing-suite
pip install -r requirements.txt
pytest -v
\end{lstlisting}
you should now see \textbf{69 passed, 2 xfailed, 1 failed}: four more
passing tests than Week 5's baseline, and the same two \texttt{xfail}
results and one \texttt{FAIL} as before, unrelated to this week's
additions. \path{demos/data_flow/} is not part of that run
(\texttt{pytest.ini} sets \texttt{testpaths = tests}); run it by path.
Its one failing test is intentional.

% ============================================================
\section{How This Connects to the Reading}
% ============================================================
\begin{keyidea}
Week 5's coverage report told you \emph{which code ran}. This week's
files ask \emph{whether the right values got where they needed to go}.
The same function, \path{sum_positive()}, appears twice: once correct,
with tests chosen pair by pair, and once broken, with a test set that
satisfies All-Defs and 100\% branch coverage and still misses the bug.
The difference between them is one DU pair.
\end{keyidea}

% ============================================================
\section{Suggested Next Steps}
% ============================================================
\begin{enumerate}
\item Before opening \path{tests/billing/test_charges.py}, list the seven
DU pairs of \path{sum_positive()} yourself and choose a minimal All-Uses
test set. Then compare with the test ids.
\item Fix \path{demos/data_flow/accumulator_bug.py} and confirm its
failing test passes, with the coverage report unchanged at 100\%.
\item Fix both functions in \path{demos/data_flow/anomalies.py} so that
\texttt{pylint} reports nothing. Add the boundary test that should
have caught the first one.
\item Work the reading's \path{first_negative()} exercises on paper,
then write the function and a parametrized test whose ids name the DU
pairs, in the style of \path{test_charges.py}.
\item The Week 5 coverage gap in \path{app/orders/order_service.py}
is still open. If you have not closed it yet, do that too.
\end{enumerate}

\end{document}
